\section{拉普拉斯变换求解微分方程}\label{sec:Laplace_ODE}

\subsection{拉普拉斯变换的改进}\label{sub:improvement_of_laplace}

由于$f\in\mathcal{E}$对函数增长速度的要求是相当低的（作为反例，$\mathrm{e}^{t^2}\notin\mathcal{E}$），我们只需要单独处理$\delta$分布，而不需要重新构建一整套分布的理论\footnote{希望深入了解分布及其拉普拉斯变换的读者，可以参考\cite{Laplace_Distribution}。\cite{distribution}讨论这个话题时，采用了与本书类似的处理方法，这实际上不是拉普拉斯变换的分布理论。}。仍考虑将拉普拉斯变换化归为傅里叶变换，即
\[\mathcal{L} f(s)=\mathcal{F} [f(t)\mathrm{e}^{-\sigma t}](\omega),s=\sigma+\mathrm{i}\omega\]
这里$\mathcal{F} [f(t)\mathrm{e}^{-\sigma t}](\omega)$实际上依赖于两个变量$\sigma,\omega$，但我们只写出$\omega$，并将$\sigma$视为一个参数\cite{distribution}。这样一来，我们就可以将拉普拉斯变换推广到$\delta$分布上（为了避免变量符号的问题，我们将变量统一地用$x$表示）：
\begin{align*}
    \left\langle \mathcal{L}\delta,\phi\right\rangle=\left\langle \mathrm{e}^{-\sigma x}\delta,\mathcal{F} \varphi\right\rangle=\left\langle \delta,\mathrm{e}^{-\sigma x}\mathcal{F} \varphi\right\rangle=\mathcal{F}\varphi(0)=\int_{-\infty}^{\infty}\varphi(t)\,dt=\left\langle \mathds{1},\varphi\right\rangle
\end{align*}
其中$\mathds{1}$表示恒为1的函数。我们知道，多项式的拉普拉斯变换是真分式函数，而如果将$\delta$分布也纳入考虑，则可以得到任意有理函数。

用同样的方法，我们还可以得到$\delta$的平移和导数的拉普拉斯变换：
\begin{align*}
    \left\langle \mathcal{L}\delta_a,\varphi\right\rangle&=\left\langle \mathrm{e}^{-\sigma x}\delta_a,\mathcal{F} \varphi\right\rangle\\
    &=\mathrm{e}^{-\sigma a}\mathcal{F}\varphi(a)\\
    &=\int_{-\infty}^{\infty}\varphi(t)\mathrm{e}^{-s a}\,dt\\
    &=\left\langle \mathrm{e}^{-s a},\varphi\right\rangle
\end{align*}
\begin{align*}
    \left\langle \mathcal{L}\left[\delta'\right],\varphi\right\rangle&=\left\langle \mathrm{e}^{-\sigma x}\delta',\mathcal{F} \varphi\right\rangle\\
    &=-\left\langle \delta,\left(\mathrm{e}^{-\sigma x}\mathcal{F} \varphi\right)'\right\rangle\\
    &=-\left\langle \delta,-\sigma \mathrm{e}^{-\sigma x}\mathcal{F} \varphi + \mathrm{e}^{-\sigma x}(\mathcal{F} \varphi)'\right\rangle\\
    &=\left\langle \delta,\sigma \mathrm{e}^{-\sigma x}\mathcal{F} \varphi + \mathrm{i}\omega \mathrm{e}^{-\sigma x}\mathcal{F} \varphi\right\rangle\\
    &=(\sigma + \mathrm{i}\omega)\mathcal{F} \varphi(0)\\
    &=\left\langle s,\varphi\right\rangle
\end{align*}
因此我们得到：
\begin{proposition}
    设$a\in\mathbb{R}$，则$\mathcal{L} \delta_a=\mathrm{e}^{-s a}$；$\mathcal{L} \left[\delta'\right]=s$，进一步可得$\mathcal{L} \left[\delta^{(n)}\right]=s^n$。
\end{proposition}

看上去，$\delta(0)=\infty,\mathcal{L} \delta'=\mathcal{L} \delta-\infty$，而$\delta_a$在$a<0$时非因果，但$\delta$应作为一个分布来理解，因此以上结果与我们之前所定义的拉普拉斯变换并不矛盾。

不过，我们可以通过另一种方式，几乎完全避开分布的问题。在本节中，我们将（单边）拉普拉斯变换的定义修改为：
\begin{definition}[拉普拉斯变换，改进版]
    设函数$f$局部可积，且$\exists a\in\mathbb{R}$，使得$f(t)\mathrm{e}^{-at}$绝对可积，则对$\text{Re}\ s\in[a,\infty)$，$f$的拉普拉斯变换定义为
    \[\mathcal{L} f(s)=\int_{0-}^{\infty}f(t)\mathrm{e}^{-st}\,dt\]
\end{definition}
这里，积分限中的“0-”的含义是：
\[\mathcal{L} f(s)=\lim_{M\to 0-}\int_{M}^{\infty}f(t)\mathrm{e}^{-st}\,dt=\int_{0-}^{\infty}f(t)\mathrm{e}^{-st}\,dt\]

在这种定义下，拉普拉斯变换的各种运算性质仍成立，但两个性质略有差别：
\begin{proposition}[改进的拉普拉斯变换的性质]
    设$f\in C^1[0,\infty),f'$局部可积，且$\exists a\in\mathbb{R}$，使得$f(t)\mathrm{e}^{-at}$绝对可积，则对$\text{Re}\ s\in[a,\infty)$，有：
    \begin{itemize}
        \item 时域微分性质：\[\mathcal{L} [f'](s)=s\mathcal{L} f(s)-f(0-)\]
        \item 初值定理：$$f(0-)=\lim_{\text{Re}\ s\to\infty}sF(s)$$
    \end{itemize}
\end{proposition}
\begin{proof}
    \begin{align*}
        \mathcal{L} [f'](s)&=\int_{0-}^{\infty}f'(t)\mathrm{e}^{-st}\,dt\\
        &=\evalat{f(t)\mathrm{e}^{-st}}{0-}{\infty}+\int_{0-}^{\infty}f(t)se^{-st}\,dt\\
        &=s\mathcal{L} f(s)-f(0-)
    \end{align*}
    对于初值定理，只要在\ref{sub:initial_and_final_value}给出的证明中使用这个时域微分性质即可。
\end{proof}

形式上，$u(0-)=\delta(0-)=0$，自然得到
\begin{align*}
    \mathcal{L} u(s)=\frac{1}{s}&\implies\mathcal{L} \delta(s)=s\mathcal{L} u(s)-u(0-)=\mathds{1}\\
    &\implies\mathcal{L} \delta'(s)=s\mathcal{L} \delta(s)-\delta(0-)=s
\end{align*}

$\delta$的各阶导数同理。这样，分布理论处理的问题仅仅在于$\delta_a(a<0)$，但我们很少用到这个分布。

可以看到，这样的定义方式帮助我们避开了$u(t)$跳变的问题和$\delta(0)=\infty$的问题，同时积分下限为“0-”排除了$f$在$(-\infty,0)$上的行为对拉普拉斯变换的影响，仅仅是拉普拉斯变换的改进，而没有对任何性质产生本质影响。另外，这样的定义方式有助于我们讨论微分方程描述的系统的起始状态，下一小节就来讨论这个问题。

\subsection{常系数线性微分方程}\label{sub:laplace_ODE}
拉普拉斯变换是处理常系数线性微分方程的有力工具。根据拉普拉斯变换的时域微分性质：
\[\mathcal{L}\left[f^{(n)}(t)\right](s)=s^n\mathcal{L} f(s)-s^{n-1}f(0-)-s^{n-2}f'(0-)-\cdots -f^{(n-1)}(0-)\]
可以将常系数线性微分方程转化为代数方程，还能够一并处理起始状态\footnote{数学上，我们使用经典的拉普拉斯变换以处理带初始状态的初值问题。初始状态、初值问题的定义，见\ref{sub:concepts_of_system}，起始条件、初始条件的定义，见\ref{sub:term_in_ODE}}。我们先提出一个一般的解法，再通过一个具体的例子来展示拉普拉斯变换的威力。

考虑常系数线性微分方程
\[a_n r^{(n)}(t)+a_{n-1}r^{(n-1)}(t)+\cdots +a_1 r'(t)+a_0 r(t)=e(t)\]
其中$f,r\in\mathcal{E} $，并且给定初始状态：
\[r(0-)=c_0,r'(0-)=c_1,\cdots,r^{(n-1)}(0-)=c_{n-1}\]
记$\mathcal{L} r(s)=R(s),\mathcal{L} e(s)=E(s)$，则对方程两边进行拉普拉斯变换，有
\begin{align*}
    &a_n\left(s^n R(s)-s^{n-1}c_0 -s^{n-2}c_1 -\cdots -c_{n-1}\right)\\
    &+a_{n-1}\left(s^{n-1} R(s)-s^{n-2}c_0 -s^{n-3}c_1 -\cdots -c_{n-2}\right)\\
    &+\cdots +a_1\left(s R(s)-c_0\right)+a_0 R(s)=E(s)
\end{align*}
整理得
\[P(s)R(s)=E(s)+Q(s)\]
其中\begin{align*}
P(s)&=a_n s^n +a_{n-1} s^{n-1}+\cdots +a_1 s+a_0\\
Q(s)&=a_n\left(s^{n-1}c_0 +s^{n-2}c_1 +\cdots +c_{n-1}\right)+a_{n-1}\left(s^{n-2}c_0 +s^{n-3}c_1 +\cdots +c_{n-2}\right)+\cdots +a_1 c_0
\end{align*}
因此
\[R(s)=\frac{E(s)}{P(s)}+\frac{Q(s)}{P(s)}=R_{zs}(s)+R_{zi}(s)\]
现在，微分方程问题就转化为求$R(s)$的拉普拉斯逆变换的问题，并且可以求出系统函数
\[H(s)=\frac{R_{zs}(s)}{E(s)}=\frac{1}{P(s)}\]
与\ref{sub:fundamental_solution}类似，常系数线性微分方程所描述的系统作为一个线性时不变系统，其零状态响应的求解也可以归结为单位脉冲响应，即基本解的求解。现在，我们找到了一种新的、更为方便的求解方法，即先由拉普拉斯变换求得系统函数在做拉普拉斯逆变换。

\begin{example}[常系数线性微分方程]
    设系统由如下微分方程描述：
    \[r''(t)+3r'(t)+2r(t)=e'(t)+3e(t)\]
    初始状态为$r(0-)=1,r'(0-)=2$。求系统函数$H(s)$及单位阶跃响应$g(t)$。\\
    解：对方程两边进行拉普拉斯变换，有（注意$e(0-)=u(0-)=0$）
    \begin{align*}
        \left(s^2 R(s)-s-2\right)+3\left(s R(s)-1\right)+2 R(s)=sE(s)+3E(s)
    \end{align*}
    整理得
    \[R(s)=\frac{(s+3)E(s)+s+5}{s^2+3s+2}\]
    因此系统函数为
    \[H(s)=\frac{R_{zs}(s)}{E(s)}=\frac{s+3}{s^2+3s+2}=\frac{2}{s+1}-\frac{1}{s+2}\]
    做拉普拉斯反变换得到（零状态）单位脉冲响应：
    \begin{align*}
        h(t)=\mathcal{L}^{-1} \left[\frac{1}{s+1}+\frac{1}{s+2}\right](t)=(2\mathrm{e}^{-t}-\mathrm{e}^{-2t})u(t)
    \end{align*}
    积分得到零状态响应：
    \begin{align*}
    r_{zs}(t)=\int_0^{t}h(x)\,dx=(1.5-2\mathrm{e}^{-t}+0.5\mathrm{e}^{-2t})u(t)
    \end{align*}
    对于零输入响应，有
    \begin{align*}
        R_{zi}(s)=\frac{s+5}{s^2+3s+2}=\frac{4}{s+1}-\frac{3}{s+2}
    \end{align*}
    反变换即得$$r_{zi}(t)=(4\mathrm{e}^{-t}-3\mathrm{e}^{-2t})u(t)$$
    实际上，零输入响应并不是因果的，但由于单边拉普拉斯变换自动舍弃了负半轴的信息，我们只能得到$t\geq 0$时的信号情况。不难看出$r_{zi}(t)=4\mathrm{e}^{-t}-3\mathrm{e}^{-2t},t<0$。

    我们一般不考虑$t<0$的情况，而直接将全响应写为
    $$r(t)=r_{zs}(t)+r_{zi}(t)=(1.5+2\mathrm{e}^{-t}-2.5\mathrm{e}^{-2t})u(t)$$

    此外，也可以用$\mathcal{L} u(s)=1/s$直接计算单位阶跃响应得全响应：
    \begin{align*}
        R(s)&=\frac{s+3}{s(s^2+3s+2)}+\frac{s+5}{s^2+3s+2}\\
        &=\frac{1.5}{s}+\frac{2}{s+1}-\frac{2.5}{s+2}
    \end{align*}
    反变换即得$r(t)=(1.5+2\mathrm{e}^{-t}-2.5\mathrm{e}^{-2t})u(t)$。
\end{example}

\subsection{动态电路的复频域分析}\label{sub:laplace_circuit}

本小节介绍动态电路的复频域分析方法。与\ref{sub:freq_response}节类似，对描述线性元件的微分方程做拉普拉斯变换，即可得到它们的复频域模型：
\begin{example}
    描述电容的方程为：
    \[\mathrm{i}(t)=C\frac{dv(t)}{dt}\]
    对其做拉普拉斯变换，得到
    \[I(s)=C\left(sV(s)-v(0-)\right)\]
    由于$v(0-)$未必为0，不能直接用阻抗$1/sC$来描述电容，需要引入电容的起始电压源模型：
    \[V(s)=\frac{1}{sC}I(s)+\frac{v(0-)}{s}\]
    \begin{center}
        \begin{circuitikz}[american, scale=0.8, transform shape]
            \draw (0,0) to[short, o-] (1,0)
            to[C=$1/sC$] (3.5,0)
            to[V=$v(0-)/s$] (6,0)
            to[short, -o] (7,0);
        \end{circuitikz}
    \end{center}
    描述电感的方程为：
    \[v(t)=L\frac{di(t)}{dt}\]
    对其做拉普拉斯变换，得到
    \[V(s)=L\left(sI(s)-\mathrm{i}(0-)\right)\]
    引入电感的起始电流源模型：
    \begin{center}
        \begin{circuitikz}[american, scale=0.8, transform shape]
            \draw (0,0) to[short, o-] (1,0)
            to[L=$sL$] (3.5,0)
            to[I=$Li(0-)$] (6,0)
            to[short, -o] (7,0);
        \end{circuitikz}
    \end{center}
    在零状态条件下，起始电压源/电流源均为0，可以直接使用阻抗模型。
\end{example}

作为一个例子，我们使用复频域分析的方法来处理一个二阶动态电路：
\begin{example}[二阶动态电路]
    设电路如下图(a)所示，起始状态为0，$t=0$时刻开关闭合，求电流$\mathrm{i}(t)$。
    \begin{center}
    \begin{tabular}{cc}
        % 左侧电路图 (时域)
        \begin{circuitikz}[american]
            \draw (0,0) to[V=$E$] (0,3)
            to[closing switch] (2.5,3)
            to[L=$L$, i=$\mathrm{i}(t)$] (5,3)
            to[C=$C$] (5,0)
            to[R=$R$] (0,0);
        \end{circuitikz}
        &
        % 右侧电路图 (s域)
        \begin{circuitikz}[american]
            \draw (0,0) to[V=$E/s$] (0,3)
            to[short] (2.5,3)
            to[L=$sL$, i=$I(s)$] (5,3)
            to[C=$1/sC$] (5,0)
            to[R=$R$] (0,0);
        \end{circuitikz}
        \\
        \text{(a) 时域电路} & \text{(b) s域等效电路}
    \end{tabular}
\end{center}
    解：采用复频域模型，如上图(b)所示。由基尔霍夫电压定律，有
    \[\frac{E}{s}=I(s)\left(R+sL+\frac{1}{sC}\right)\]
    因此
    \[I(s)=\frac{E}{s\left(R+sL+\frac{1}{sC}\right)}=\frac{E}{Ls^2+Rs+\frac{1}{C}}\]
    引入参数
    \[\alpha=\frac{R}{2L},\omega_0=\frac{1}{\sqrt{LC}}\]
    则
    \[I(s)=\frac{E/L}{s^2+2\alpha s+\omega_0^2}=\frac{E/L}{(s+\alpha)^2+\omega_0^2-\alpha^2}\]
    \ding{172}设$\alpha>\omega_0$，即$R^2-4L/C>0$，则
    \[I(s)=\frac{E/L}{(s-r_1)(s-r_2)}\]
    其中$$r_1=\frac{-R+\sqrt{R^2-4L/C}}{2L}=-\alpha+\beta,r_2=\frac{-R-\sqrt{R^2-4L/C}}{2L}=-\alpha-\beta,\beta=\frac{\sqrt{R^2-4L/C}}{2L}$$
    做拉普拉斯逆变换，得到
    \[\mathrm{i}(t)=\frac{E}{L(r_1-r_2)}\left(\mathrm{e}^{r_1 t}-\mathrm{e}^{r_2 t}\right)u(t)=\frac{E}{2L\beta}\left[\mathrm{e}^{(-\alpha+\beta) t}-\mathrm{e}^{(-\alpha-\beta) t}\right]u(t)\]
    系统处于\textbf{过阻尼状态}，电流为两项指数衰减的叠加，缓慢无振荡地回到平衡位置。\\
    \ding{173}设$\alpha=\omega_0$，即$R^2-4L/C=0$，则
    \[I(s)=\frac{E}{L}\cdot\frac{1}{(s+\alpha)^2}\]
    做拉普拉斯逆变换，得到
    \[\mathrm{i}(t)=\frac{E}{L}te^{-\alpha t}u(t)\]
    系统处于\textbf{临界阻尼状态}，电流快速无振荡地回到平衡位置。\\
    \ding{174}设$\alpha<\omega_0$，即$R^2-4L/C<0$，则
    \[I(s)=\frac{E}{L}\cdot\frac{1}{(s+\alpha)^2+\beta^2}\]
    做拉普拉斯逆变换，得到
    \[\mathrm{i}(t)=\frac{E}{L\beta}\mathrm{e}^{-\alpha t}\sin(\beta t)u(t)\]
    系统处于\textbf{欠阻尼状态}，电流以衰减振荡的形式缓慢回到平衡位置。
    \begin{figure}[H]
        \centering
        \includegraphics[width=0.8\textwidth]{damp}
        \caption{不同阻尼状态下的电流响应}
    \end{figure}
\end{example}

\subsection{常系数线性微分方程组的拉普拉斯变换解法}\label{sub:laplace_ode_system}

拉普拉斯变换同样能够处理常系数线性微分方程组，并且能够一并处理初始条件。我们首先对矩阵值函数定义改进版的拉普拉斯变换：
\begin{definition}[矩阵值函数的拉普拉斯变换]
    设矩阵值函数$A(t)$的每个分量局部可积，并且对于每个分量$a_{ij}(t)$，$\exists \alpha_{ij}\in\mathbb{R}$，使得$a_{ij}(t)\mathrm{e}^{-\alpha_{ij} t}$绝对可积，则在所有分量的公共收敛域
    $$\text{Re}\ s\geq \max\{\alpha_{ij}|1\leq \mathrm{i}\leq m,1\leq j\leq n\}$$上，矩阵值函数$A(t)$的拉普拉斯变换是每个分量的拉普拉斯变换构成的矩阵：
    \[\mathcal{L} A(s)=\int_{0-}^{\infty}A(t)\mathrm{e}^{-st}\,dt\]
\end{definition}

考虑如下常系数线性微分方程组：
\begin{align*}
    \frac{d\mathbf{y}}{dt}=A\mathbf{y}+\mathbf{f}(t),\mathbf{y}(0-)=\mathbf{y}_0
\end{align*}
其中$A\in\mathbb{R} ^{n\times n},\mathbf{y},\mathbf{f}\in\mathcal{E} ^n$，并且$\mathbf{y}_0\in\mathbb{R} ^n$。对方程两边进行拉普拉斯变换，有
\begin{align*}
    s\mathbf{Y}(s)-\mathbf{y}_0=A\mathbf{Y}(s)+\mathbf{F}(s)
\end{align*}
其中$\mathbf{Y}(s)$和$\mathbf{F}(s)$分别是$\mathbf{y}(t)$和$\mathbf{f}(t)$的拉普拉斯变换。因此
\begin{align*}
    \mathbf{Y}(s)=(sI-A)^{-1}(\mathbf{y}_0+\mathbf{F}(s))
\end{align*}
两边同时求拉普拉斯逆变换，即可得到方程的解：\begin{align*}
    \mathbf{y}(t)=\mathcal{L} ^{-1} \left[(sI-A)^{-1}\mathbf{y}_0\right](t)+\mathcal{L} ^{-1} \left[(sI-A)^{-1}\mathbf{F}(s)\right](t)
\end{align*}
根据拉普拉斯变换的卷积定理，令$\mathcal{L} ^{-1} \left[(sI-A)^{-1}\right](t)=\phi(t)u(t)$，则解为\begin{align*}
    \mathbf{y}(t)&=\phi(t)u(t)\mathbf{y}_0+(\phi(t)u(t))*\mathbf{f}(t)\\
    &=\phi(t)u(t)\mathbf{y}_0+\int_0^{t}\phi(t-\tau)\mathbf{f}(\tau)\,d\tau,t>0
\end{align*}
可以看出，$\phi(t)$正是方程的基本解矩阵，即
$$\mathcal{L}^{-1} \left[(sI-A)^{-1}\right](t)=\phi(t)u(t)=\mathrm{e}^{At}u(t)$$
这样，我们就得到了第三章中提到的矩阵指数函数$\mathrm{e}^{At}$的拉普拉斯变换求法。

\begin{definition}[特征矩阵/预解矩阵]
    设系统的状态矢量为$\mathbf{q}(t)$，由以下状态变量方程方程描述：
    $$\frac{d\mathbf{q}}{dt}=A\mathbf{q}+\mathbf{f}(t)$$
    求解该方程的关键是求矩阵$(sI-A)^{-1}=\mathcal{L}\{\mathrm{e}^{At}\}$。我们将矩阵$(sI-A)^{-1}$称为系统的\textbf{特征矩阵}(characteristic matrix)或\textbf{预解矩阵}(pre-solution matrix)，记为$\Phi(s)$。
\end{definition}

称$(sI-A)^{-1}$为特征矩阵，是因为它的行列式的极点就是矩阵$A$的特征值\footnote{$(sI-A)$的每个分量都是分式函数，因此其逆矩阵的每个分量也是分式函数，行列式只会有极点而不会有其它形式的奇点}，在状态变量的可逆线性变换下，特征值是不变的，因此我们也将系统矩阵$A$的特征值称为\textbf{系统的特征值}。

\begin{example}
    设系统由如下状态变量方程描述：
    \begin{align*}
        \frac{d}{dt}\begin{bmatrix}
            q_1(t)\\
            q_2(t)
        \end{bmatrix}=\begin{bmatrix}
            0 & 1\\
            -2 & -3
        \end{bmatrix}\begin{bmatrix}
            q_1(t)\\
            q_2(t)
        \end{bmatrix}+\begin{bmatrix}
            0\\
            t
        \end{bmatrix},\begin{bmatrix}
            q_1(0-)\\
            q_2(0-)
        \end{bmatrix}=\begin{bmatrix}
            1\\
            2
        \end{bmatrix}
    \end{align*}
    则特征矩阵为\begin{align*}
        \Phi(s)=(sI-A)^{-1}=\begin{bmatrix}
            s & -1\\
            2 & s+3
        \end{bmatrix}^{-1}=\frac{1}{s^2+3s+2}\begin{bmatrix}
            s+3 & 1\\
            -2 & s
        \end{bmatrix}
    \end{align*}
    可以看出，系统的特征值为$-1,-2$。对$\Phi(s)$求拉普拉斯逆变换，即可得到系统的基本解矩阵$\mathrm{e}^{At}$：\begin{align*}
        \mathrm{e}^{At}u(t)&=\mathcal{L} ^{-1} \left[\Phi(s)\right](t)\\
        &=\mathcal{L} ^{-1} \left[\frac{1}{s^2+3s+2}\begin{bmatrix}
            s+3 & 1\\
            -2 & s
        \end{bmatrix}\right](t)\\
        &=\begin{bmatrix}
            2\mathrm{e}^{-t}-\mathrm{e}^{-2t} & \mathrm{e}^{-t}-\mathrm{e}^{-2t}\\
            -2\mathrm{e}^{-t}+2\mathrm{e}^{-2t} & -2\mathrm{e}^{-t}+3\mathrm{e}^{-2t}
        \end{bmatrix}u(t)
    \end{align*}
    状态变量方程的解为\begin{align*}
        \begin{bmatrix}
            q_1(t)\\
            q_2(t)
        \end{bmatrix}&=\mathrm{e}^{At}\begin{bmatrix}
            1\\
            2
        \end{bmatrix}+\int_0^{t}\mathrm{e}^{A(t-\tau)}\begin{bmatrix}
            0\\
            \tau
        \end{bmatrix}\,d\tau\\
        &=\begin{bmatrix}
            4\mathrm{e}^{-t}-5\mathrm{e}^{-2t}\\
            -6\mathrm{e}^{-t}+8\mathrm{e}^{-2t}
        \end{bmatrix}+\int_0^{t}\begin{bmatrix}
            \mathrm{e}^{-(t-\tau)}-\mathrm{e}^{-2(t-\tau)}\\
            -2\mathrm{e}^{-(t-\tau)}+3\mathrm{e}^{-2(t-\tau)}
        \end{bmatrix}\tau\,d\tau\\
        &=\begin{bmatrix}
            5\mathrm{e}^{-t}-\frac{11}{2}\mathrm{e}^{-2t}+te^{-t}-\frac{7}{2}te^{-2t}-\frac{1}{2}\\
            -4\mathrm{e}^{-t}+\frac{13}{2}\mathrm{e}^{-2t}-2te^{-t}+\frac{9}{2}te^{-2t}-\frac{3}{2}
        \end{bmatrix}
    \end{align*}
\end{example}

考虑一个$q$输入、$p$输出的线性时不变系统，其状态变量方程为
\begin{align*}
    \frac{d\mathbf{q}}{dt}=&A\mathbf{q}+B\mathbf{e},\\
    \mathbf{r}=&C\mathbf{q}+D\mathbf{e}
\end{align*}
其中$A$为$n$阶矩阵，$n$为系统的状态变量个数；$B,C,D$分别为$n\times q$、$p\times n$和$p\times q$的矩阵。

要求其系统函数，使用输出与输入的复频域形式之比的定义较为方便。对状态变量方程两边进行拉普拉斯变换，有
\begin{align*}
    s\mathbf{Q}(s)-\mathbf{q}(0-)=&A\mathbf{Q}(s)+B\mathbf{E}(s),\\
    \mathbf{R}(s)=&C\mathbf{Q}(s)+D\mathbf{E}(s)
\end{align*}
由前一方程，
\begin{align*}
    \mathbf{Q}(s)=(sI-A)^{-1}\mathbf{q}(0-)+(sI-A)^{-1}B\mathbf{E}(s)=\Phi(s)\mathbf{q}(0-)+\Phi(s)B\mathbf{E}(s)
\end{align*}
代入后一方程，得到
\begin{align*}
    \mathbf{R}(s)=&C\left[\Phi(s)\mathbf{q}(0-)+\Phi(s)B\mathbf{E}(s)\right]+D\mathbf{E}(s)\\
    =&C\Phi(s)\mathbf{q}(0-)+\left[C\Phi(s)B+D\right]\mathbf{E}(s)
\end{align*}
因此系统函数为
\begin{align*}
    H(s)=\frac{R_{zs}(s)}{E(s)}=C\Phi(s)B+D
\end{align*}
这是一个$p\times q$的矩阵值函数。特别地，对于单输入单输出系统，系统函数为一个分式函数。

利用系统函数，可以得到系统的可控性、可观性的其他判断方法，但由于这个方法比较复杂，这里仅仅给出它的一个简化版本：
\begin{claim}
    设单输入单输出系统的状态变量方程为
    \begin{align*}
        \frac{d\mathbf{q}}{dt}=&A\mathbf{q}+B e,r=C\mathbf{q}
    \end{align*}
    则系统的系统函数为
    \begin{align*}
        H(s)=\frac{R_{zs}(s)}{E(s)}=C\Phi(s)B
    \end{align*}
    如果在系统函数的计算过程中，没有出现零极点的抵消，即系统函数的分子和分母没有公共因子，则系统是可控且可观的，否则系统不完全可控或不完全可观。
\end{claim}
